\documentclass[twoside]{article}
\usepackage{graphicx}
\usepackage{polski}
\usepackage{fancyhdr}
\usepackage{lastpage}
\usepackage{listings}
\usepackage{soul}
\usepackage{color}
\usepackage{multirow}
\usepackage{subcaption}

\title{AiSD Ex. 1 (Pon.)}
\author{\hl{Mateusz Kamieniecki}}
\date{\today}

\makeatletter

\begin{document}

\pagestyle{fancy}
\fancyhead{}
\fancyhead[L]{\@title}
\fancyhead[R]{\@author}
\fancyhead[C]{\@date}

\cfoot{\thepage\ / \pageref{LastPage}}


\newpage

\section{Wprowadzenie}
Pierwsze ćwiczenie, którym się zajmiemy, będzie związane z problemem sortowania danych.
Dla ułatwienia dalszych rozważań przyjmijmy, że mówiąc o~sortowaniu kolekcji liczb (lub łańcuchów znaków), 
będziemy rozważali porządek leksykograficzny.

Zarówno w życiu codziennym jak i w praktyce inżynierskiej dany problem może mieć wiele rozwiązań.
Nie inaczej jest z zadaniem sortowania, dla którego powstało wiele różnych algorytmów.
Wybór odpowiedniego algorytmu może być podyktowany szeregiem czynników, do których można zaliczyć 
m.in.: łatwość samodzielnej implementacji, złożoność obliczeniową (czasową, pamięciową), 
dostępność zasobów.
W przypadku rozwiązywania problemów nieszablonowych, dla których nie istnieje gotowa instrukcja
postępowania, poszukiwanie dobrego rozwiązania jest kompromisem pomiędzy satysfakcją 
z~uzyskanego wyniku (złożonością algorytmu), a czasem poświęconym na jego implementację.
Umiejętność świadomego wyważenia tych proporcji jest kompetencją inżynierską poszukiwaną 
na rynku pracy.

\subsection{Złożoność algorytmów}
Ocena czasu pracy algorytmu jak i zużywanych przez niego zasobów pamięciowych może być wykonana bardzo dokładnie
np. poprzez zliczenie wszystkich instrukcji, które zostaną wykonane w trakcie jego działania.

Częściej jednak mówimy o klasach złożoności algorytmów przybliżając ich precyzyjnie wyliczone wartości
asymptotycznymi złożonościami, co wynika z~faktu, że najczęściej precyzyjnie wyliczone wartości 
są funkcjami rozmiaru danych $n$ (np.: $f(n) = an^2 + bn + c$, gdzie $a > 0$), a przy bardzo dużym rozmiarze danych
$n$ przyjmuje się, że składniki niższych rzędów i stały współczynnik przy składniku o najwyższym rzędzie stają się nieistotne (są pomijane w rozważaniach).

Do najpopularniejszych notacji zalicza się:
\begin{itemize}
    \item notację $O$ -- do określenia asymptotycznego ograniczenia górnego;
    \item notację $\Omega$ -- do określenia asymptotycznego ograniczenia dolnego;
    \item notację $\Theta$ -- do określenia asymptotycznego dokładnego oszacowania.
\end{itemize}

Na podstawie tych rozważań można napisać, że funkcja $f(n)$ ($f(n) = an^2 + bn + c$, gdzie $a > 0$), określająca
pesymistyczny czas działania algorytmu dla danych o rozmiarze $n$, jest $O(n^2)$ ($f(n)=O(n^2)$).

Biorąc to wszystko pod uwagę, zastanówmy się, jaka może być złożoność wyrażona w notacji $O$ dla poniższych 
fragmentów kodu (\ref{lst:big_o1}, \ref{lst:big_o2}), przy założeniu, że $numsA$ i $numsB$ składają się 
odpowiednio z $a$ i $b$ elementów.

\begin{minipage}[t]{.45\textwidth}
\begin{lstlisting}[caption={Złożoność $O(a+b)$.}, frame=tblr, language=Java, label={lst:big_o1}]
for(int a: numsA){
    .println(a);
}

for(int b: numsB){
    .println(b);
}
\end{lstlisting}
\end{minipage}
\hfill
\begin{minipage}[t]{.45\textwidth}
\begin{lstlisting}[caption={Złożoność $O(a*b)$.}, frame=tblr, language=Java, label={lst:big_o2}]
for(int a: numsA){
    for(int b: numsB){
        .println(a + b);
    }
}
\end{lstlisting}
\end{minipage}

Ma to sens, prawda?
To spróbujmy jeszcze dwa razy (listingi: \ref{lst:big_o3}, \ref{lst:big_o4}).

W przypadku listingu \ref{lst:big_o3} wykonujemy dwie niezależne pętle, a w każdej z nich kręcimy się $n$ razy,
dlatego złożoność z $O(2n)$ redukujemy (pomijając stałe) do $O(n)$.

\begin{lstlisting}[caption={Złożoność $O(n)$.}, frame=tblr, language=Java, label={lst:big_o3}]
void countProductAndSum(int[] nums){
    int product = 1;
    int sum = 0;
    int n = nums.length;
    
    for(int i = 0; i < n; i++){
        product *= nums[i];
    }
    
    for(int i = 0; i < n; i++){
        sum += nums[i];
    }
    
    System.out.println(
        format("Product: %d | Sum: %d", product, sum));
}
\end{lstlisting}

Ostatni przykład (listing \ref{lst:big_o4}) pokazuje nam dwie zagnieżdżone pętle -- dla każdego kroku ($n$ kroków)
z pętli zewnętrznej wykonujemy n-krokową pętlę wewnętrzną, więc uzyskujemy złożoność $O(n^2)$.

\begin{lstlisting}[caption={Złożoność $O(n^2)$.}, frame=tblr, language=Java, label={lst:big_o4}]
void printEachWithEach(int[] nums){
    int n = nums.length;
    
    for(int i = 0; i < n; i++){
        for(int j = 0; j < n; j++){
            System.out.println(
                format("%d, %d", nums[i], nums[j]));
        }
    }
}
\end{lstlisting}

\subsection{Algorytmy sortowania}

Praca na uporządkowanej (posortowanej) kolekcji elementów może być celem samym w sobie, ale może być również
pewnym etapem przejściowym prowadzącym do celu. Niezależnie od powodu, dla którego chcemy otrzymać uporządkowaną
kolekcję, warto jest znać różnice (zalety i wady) stosowania różnych algorytmów służących do tego celu.

Algorytmy różnią się pod względem m.in.: złożoności obliczeniowej (czasowej, pamięciowej), trudności implementacyjnej
(gdy nie korzystamy z bibliotek gotowych implementacji), stabilnością sortowania, pracą w miejscu 
(łac. \textit{in-situ}), szczególnym zachowaniem w przypadku określonego układu lub rozmiaru danych wejściowych.

Przykładowo, algorytm sortowania przez wybór (ang. \textit{selection sort}) ma praktycznie taką samą złożoność
czasową $O(n^2)$ niezależnie od układu danych wejściowych. A algorytm sortowania przez wstawianie (ang. 
\textit{insertion sort}) dla małych rozmiarów danych wejściowych lub danych prawie uporządkowanych 
(wstępnie posortowanych) działa znacznie szybciej (nawet $\Theta(n)$).

Omówienie zarówno podstawowych jak i bardziej specyficznych algorytmów sortowania można znaleźć we 
\textit{Wprowadzeniu do algorytmów} Cormena i in., dlatego te treści nie będą tutaj szczegółowo powielane.
Warto jednak podkreślić kilka szczególnych faktów, które będą istotne przy wykonaniu pierwszego
ćwiczenia laboratoryjnego.

Algorytm sortowania przez kopcowanie (ang. \textit{heap sort}) swoje działanie opiera na właściwościach 
tablicowej struktury danych jaką jest \textbf{kopiec binarny} -- wygodnie jest ją rozpatrywać teoretycznie jako drzewo binarne.
Najczęściej mówimy o \textbf{maksymalnym kopcu}, co oznacza, że w jego korzeniu znajduje się największa wartość, a 
relacja między rodzicem a potomkami (czyli \textbf{własność kopca}) określa, że rodzic nie może być mniejszy
od swoich dzieci.
Kopiec ten jest również \textbf{kopcem zupełnym} (drzewem), a tym samym liście w tym kopcu mogą występować tylko
na dwóch ostatnich poziomach.
Spełnienie \textbf{własności kształtu} oznacza, że nowi potomkowie dołączają do drzewa na pierwszą wolną pozycję.

Algorytm szybkiego sortowania (ang. \textit{quick sort}) jest jednym z częściej wykorzystywanych algorytmów sortujących
(wraz ze swoimi ulepszeniami i wariantami implementacyjnymi). Nie jest on jednak algorytmem, który jest najlepszym
wyborem do każdego zadania, choć zdecydowanie korzystnie wypada w użyciu dla typowych, ,,codziennych'' sytuacji.

Algorytm ten występuje w dwóch popularnych wariantach zaproponowanych przez Sir Tony'ego Hoare'a i Nico Lomuto.
Jest to algorytm z grupy ,,dziel i zwyciężaj'', co oznacza (jak wynika z nazwy), że składa się z \st{dwóch} trzech etapów: dziel, zwyciężaj, połącz. A przytoczone warianty różnią się właśnie wykorzystaną metodą podziału.

Najwygodniej jest (tylko dla celów poglądowych, ale \textbf{nigdy dla docelowej implementacji}) przedstawić algorytm
szybkiego sortowania przy pomocy rekurencji. Nie jest to jednak zalecany sposób wykonywania implementacji ze względu
na koszt (i ryzyko błędu) działania takiego kodu.

\begin{lstlisting}[caption={QuickSort (met. Lomuto) -- poglądowa wersja rekurencyjna}, frame=tblr, language=Java, label={lst:qsort_lomuto_rec}]
void quickSort(double[] nums, int left, int right) {
    if (left < right) {

        int pivotId = partition(nums, left, right);

        quickSort(nums, left, pivotId - 1);
        quickSort(nums, pivotId + 1, right);
    }
}
\end{lstlisting}

\begin{lstlisting}[caption={QuickSort (met. Hoare'a) -- poglądowa wersja rekurencyjna}, frame=tblr, language=Java, label={lst:qsort_hoare_rec}]
void quickSort(double[] nums, int left, int right) {
    if (left < right) {

        int pivotId = partition(nums, left, right);

        quickSort(nums, left, pivotId);
        quickSort(nums, pivotId + 1, right);
    }
}
\end{lstlisting}

Listingi \ref{lst:qsort_lomuto_rec}. i \ref{lst:qsort_hoare_rec}. pokazują przykładową rekurencyjną implementację
algorytmu szybkiego sortowania przy założeniu, że metoda \verb|partition(...)| wykonuje podział odpowiednio 
podejściem Lomuto lub Hoare'a.

Z pokazanego kodu wyraźnie widać (przy założeniu, że ten kod jest prawidłowy), że najistotniejszym elementem 
w algorytmie szybkiego sortowania jest właśnie procedura podziału zwracająca indeks elementu osiowego (rozdzielającego).
Celem wykonania podziału jest przeniesienie elementów w taki sposób, żeby po lewej stronie były elementy mniejsze 
(bądź równe) wartości rozdzielającej, a po prawej stronie -- elementy większe.

Różnica w podejściu Hoare'a i Lomuto związana jest z tym, który element zostaje wybrany jako element rozdzielający, 
i w którym miejscu (po podziale) ostatecznie znajduje się ten element.

Algorytmem sortującym szczególnego przeznaczenia (nie dla każdych danych) jest algorytm sortowania przez zliczanie
(ang. \textit{counting sort}). Ten algorytm dla zestawu danych (powtarzających się liczb z góry ograniczonego zakresu znacznie mniejszego niż rozmiar danych) może działać ze złożonością $O(n)$. Trzeba jednak ciągle mieć na uwadze, że
nie jest to algorytm ogólnego przeznaczenia.
A najistotniejszym novum -- względem pozostałych algorytmów -- jest wykorzystanie dodatkowej pamięci na stworzenie \verb|nOfCounters| liczników zliczających wystąpienie poszczególnych wartości przypisanych do danego kubełka.

\begin{lstlisting}[caption={Counting sort}, frame=tblr, language=Java, label={lst:counting_sort}]
void countingSort(int[] nums, int maxVal) {
    // ... validateInput

    int n = nums.length;
    int nOfCounters = maxVal + 1;

    int[] counters = new int[nOfCounters];

    for (int i = 0; i < n; i++) {
        counters[nums[i]]++;
    }

    int j = 0;
    for (int i = 0; i < nOfCounters; i++) {
        while (counters[i] > 0) {
            nums[j++] = i;
            counters[i]--;
        }
    }
}
\end{lstlisting}

\section{Zadania}

\subsection{Zadanie 1.}
Proszę, wykorzystując notację wielkiego $O$, określić asymptotyczną złożoność czasową trzech poniższych fragmentów kodu.

\begin{table}[h]
\begin{tabular}{|l|l|l|}
\hline
Lp. & Nr przykładu & Złożoność O(?) \\ \hline
1.  & A            & O(n^2)               \\ \hline
2.  & B            & O(numsA.length * numsB.length)               \\ \hline
3.  & C            & O(n)               \\ \hline
\end{tabular}
\end{table}

\begin{lstlisting}[caption={Zad. 1 (A)}, frame=tblr, language=Java]
void printHalfOfPairs(int[] nums) {
    int n = nums.length; // may be null!

    for (int i = 0; i < n; i++) {
        for (int j = i + 1; j < n; j++) {
            System.out.println(
                format("(%d, %d)", nums[i], nums[j]));
        }
    }
}
\end{lstlisting}

\begin{lstlisting}[caption={Zad. 1 (B)}, frame=tblr, language=Java]
void printSmallerPairs(int[] numsA, int[] numsB) {
    // may be nulls!
    for (int i = 0; i < numsA.length; i++) {
        for (int j = 0; j < numsB.length; j++) {
            if (numsA[i] < numsB[j]) {
                System.out.println(
                 format("(%d, %d)", numsA[i], numsB[j]));
            }
        }
    }
}
\end{lstlisting}

\begin{lstlisting}[caption={Zad. 1 (C)}, frame=tblr, language=Java]
void reverse(int[] nums) {
    int n = nums.length; // may be null!
    int rightId, leftVal;

    for (int i = 0; i < n / 2; i++) {
        rightId = n - i - 1;
        leftVal = nums[i];
        nums[i] = nums[rightId];
        nums[rightId] = leftVal;
    }
}
\end{lstlisting}

\subsection{Zadanie 2.}
Proszę określić (\textbf{T}ak lub \textbf{N}ie), który z poniższych zapisów jest tożsamy asymptotycznej złożoności $O(n)$ (wiadomo, że dla dużej wartości $n$). Proszę podać \underline{krótkie} uzasadnienie (najlepiej jako równoważnik zdania).

\begin{table}[h]
\begin{tabular}{|l|l|l|l|}
\hline
Lp. & Przykład                      & Czy jest O(n)? & Uzasadnienie \\ \hline
1.  & O(n+k), gdzie k \textless n/3 & Tak               & O(n+(1/3)n) = O((4/3)n) = O(n)              \\ \hline
2.  & O(2n)                         & Tak               & Stała ominięta             \\ \hline
3.  & O(n+m)                        & Nie               & n i m to nie stałe             \\ \hline
4.  & O(n+log(n))                   & Tak               & n rośnie szybciej niż log(n)             \\ \hline
\end{tabular}
\end{table}

\subsection{Zadanie 3.}
W kodzie źródłowym zostały umieszczone dwie klasy (\verb|QuickSortIterativeHoare|, \verb|QuickSortIterativeLomuto|) 
dla implementacji algorytmu szybkiego sortowania z wykorzystaniem metod podziału Hoare'a i Lomuto.
Proszę wybrać jedną z tych metod podziału, a następnie proszę zaimplementować trzy klasy
teoretycznie usprawniające działanie algorytmu szybkiego sortowania:

\begin{itemize}
    \item \verb|QuickSortIterativeMedian3| -- wybór elementu rozdzielającego (pivot) jako mediany z trzech wartości;
    \item \verb|QuickSortIterativeRandom| -- wybór elementu rozdzielającego (pivot) jako elementu losowo wybranego;
    \item \verb|QuickSortIterativeWithInSort| -- wykorzystanie sortowania przez wstawianie dla małych ($n < 20$) podproblemów.
\end{itemize}

\subsection{Zadanie 4.}
Proszę, wykorzystując otrzymany kod uzupełniony o własne implementacje, wykonać wykresy (nie, nie trzeba pisać kodu tworzącego wykres -- arkusz kalkulacyjny jest wystarczająco dobrym narzędziem)
pokazujący, jak zmienia się czas wykonania algorytmu w zależności od rozmiaru danych wejściowych.
Proszę wykonać wykresy w skali logarytmicznej na osi OY.
Na podstawie otrzymanych wykresów, proszę uzupełnić poniższą tabelę.

Do tego celu można wykorzystać pakiet testów \verb|performance.table|.
Zadanie dla chętnych: zamiast wspomnianego pakietu, proszę wykorzystać pakiet \verb|performance.charts|. 
Trzeba w nim wprowadzić trochę zmian, aby nie było konieczności tworzenia kilkudziesięciu klas, ale dzięki temu
wykres będzie bardziej ,,estetyczny'' (dokładniejszy).

\begin{table}[h]
    \begin{tabular}{|l|l|l|}
    \hline
    Lp. & Kategoria & \begin{tabular}[c]{@{}l@{}}Algorytm\\ (lub algorytmy)\end{tabular} \\ \hline
    1. & Najlepszy algorytm dla danych losowych &  \\ \hline
    2. & Najlepszy algorytm dla danych rosnących &  \\ \hline
    3. & Najlepszy algorytm dla danych malejących &  \\ \hline \hline
    4. & Najgorszy algorytm dla danych losowych &  \\ \hline
    5. & Najgorszy algorytm dla danych rosnących &  \\ \hline
    6. & Najgorszy algorytm dla danych malejących &  \\ \hline
    \end{tabular}
\end{table}

\begin{figure}[!htb]
    \centering
    
    \begin{subfigure}[b]{\textwidth}
        \centering
        \includegraphics[height=4.5cm]{example-image-a}
        \caption{Czas sortowania dla danych losowych.}
        \label{fig:sorters_random_data}
    \end{subfigure}
    
    \begin{subfigure}[b]{\textwidth}
        \centering
        \includegraphics[height=4.5cm]{example-image-b}
        \caption{Czas sortowania dla danych rosnących.}
        \label{fig:sorters_ascending_data}
    \end{subfigure}
    
    \begin{subfigure}[b]{\textwidth}
        \centering
        \includegraphics[height=4.5cm]{example-image-c}
        \caption{Czas sortowania dla danych malejących.}
        \label{fig:sorters_descending_data}
    \end{subfigure}
    
    \caption{Wykresy pokazują, jak zmienia się czas sortowania w zależności od rozmiaru danych dla algorytmów: sortowania przez wstawianie, sortowania przez wybór, sortowania szybkiego (różne warianty), sortowania przez scalanie, sortowania przez kopcowanie, sortowania referencyjnego.}
\end{figure}

\newpage
\subsection{Zadanie 5.}

Proszę, wykorzystując otrzymany kod uzupełniony (i poprawiony) o własne implementacje, uzupełnić klasę testów jednostkowych \verb|MergeSortStabilityTest| o metody pokazujące, czy sortowanie przez
scalanie jest algorytmem sortującym stabilnie.
Przykładowy (wybranych z samodzielnie zaimplementowanych) test proszę umieścić w listingu \ref{lst:mergeSortStable}.

\begin{lstlisting}[caption={Przykładowa metoda testowa dla MergeSort pokazująca, że algorytm jest stabilny.},
frame=tblr, language=Java, label=lst:mergeSortStable]
// TODO: placeholder for test method
\end{lstlisting}


\newpage
\subsection{Zadanie 6.}

Proszę, wykorzystując otrzymany kod uzupełniony o własne implementacje, uzupełnić poniższą tabelę 
pokazującą, jaki wpływ na czas sortowania ma typ (prosty, złożony) sortowanych danych.
W wybranych klasach konieczne jest zaimplementowanie dodatkowego interfejsu \verb|Sorting| 
(oraz dodatkowych klas testów).
Proszę \underline{krótko} skomentować otrzymane wartości.

%---------------------------------------------------------
% Please, use this tool: https://www.tablesgenerator.com/#
%---------------------------------------------------------
\begin{table}[h]
\begin{tabular}{|l|l|ll|ll|ll|}
\hline
\multirow{2}{*}{Lp.} & \multirow{2}{*}{Rozmiar} & \multicolumn{2}{c|}{InsertionSort} & \multicolumn{2}{l|}{ReferenceSort} & \multicolumn{2}{c|}{SelectionSort} \\ \cline{3-8} 
 &  & \multicolumn{1}{l|}{Integer} & int & \multicolumn{1}{c|}{Integer} & \multicolumn{1}{c|}{int} & \multicolumn{1}{l|}{Integer} & int \\ \hline
1. & 1024 & \multicolumn{1}{l|}{} &  & \multicolumn{1}{l|}{} &  & \multicolumn{1}{l|}{} &  \\ \hline
2. & 2048 & \multicolumn{1}{l|}{} &  & \multicolumn{1}{l|}{} &  & \multicolumn{1}{l|}{} &  \\ \hline
3. & 4096 & \multicolumn{1}{l|}{} &  & \multicolumn{1}{l|}{} &  & \multicolumn{1}{l|}{} &  \\ \hline
4. & 8192 & \multicolumn{1}{l|}{} &  & \multicolumn{1}{l|}{} &  & \multicolumn{1}{l|}{} &  \\ \hline
5. & 16384 & \multicolumn{1}{l|}{} &  & \multicolumn{1}{l|}{} &  & \multicolumn{1}{l|}{} &  \\ \hline
6. & 32768 & \multicolumn{1}{l|}{} &  & \multicolumn{1}{l|}{} &  & \multicolumn{1}{l|}{} &  \\ \hline
7. & 65536 & \multicolumn{1}{l|}{} &  & \multicolumn{1}{l|}{} &  & \multicolumn{1}{l|}{} &  \\ \hline
8. & 131072 & \multicolumn{1}{l|}{} &  & \multicolumn{1}{l|}{} &  & \multicolumn{1}{l|}{} &  \\ \hline
9. & 262144 & \multicolumn{1}{l|}{} &  & \multicolumn{1}{l|}{} &  & \multicolumn{1}{l|}{} &  \\ \hline
10. & 524288 & \multicolumn{1}{l|}{} &  & \multicolumn{1}{l|}{} &  & \multicolumn{1}{l|}{} &  \\ \hline
\end{tabular}
\end{table}

Komentarz (a.k.a. wnioski): \textbf{\hl{TODO -----------------------------------------}}.

\newpage
\subsection{Zadanie 7.}

Proszę, wykorzystując otrzymany kod, ustalić, jaka jest optymalna wartość podproblemów podczas
wykonywania algorytmu szybkiego sortowania na załączenie sortowania przez wstawianie.
Przeprowadzoną analizę proszę zobrazować na wykresie.

\begin{figure}[h]
    \centering
    \includegraphics[width=\linewidth]{example-image-a}
    \caption{Wykres pokazuje, jak zmienia się czas sortowania tablicy losowo uporządkowanych wartości algorytmem szybkiego sortowania (z dołączonym algorytmem sortowania przez wstawianie), w zależności od wielkości podproblemu, dla którego dołączany jest ten algorytm (sortowania przez wstawianie).}
    \label{fig:qsort_and_insort_optimal_value}
\end{figure}

Komentarz (a.k.a. wnioski): \textbf{\hl{TODO -----------------------------------------}}.

\newpage
\subsection{Zadanie 8.}

Proszę samodzielnie przygotować pakiet testów wydajnościowych dla algorytmu sortowania przez zliczanie.
Testy powinny pomóc w przygotowaniu wykresu pokazującego, jak zmienia się czas sortowania danych, 
gdy zestaw danych jest zestawem ,,przypadkowym'' (czyli bez powtórzeń), oraz w przypadku, gdy zestaw danych
jest zestawem preferowanym przez algorytm sortowania przez zliczanie (z wieloma powtórzeniami).

Miejsce na wykres oraz komentarz również należy stworzyć sobie samodzielnie w tej sekcji.

\newpage
\subsection{Zadanie 9.}

Proszę, korzystając z kodu źródłowego Javy oraz dokumentacji (Javadoc), podać, kto jest autorem
wykorzystywanego algorytmu sortowania dla tablicy liczb całkowitych (\verb|int[]|).
Proszę również określić, czym szczególnym charakteryzuje się ta implementacja.

Proszę na dowód załączyć zrzut ekranu pokazujący istotny fragment dokumentacji (Javadoc).

Odpowiedź: \textbf{\hl{TODO -----------------------------------------------------------}}.

\begin{figure}[h]
    \centering
    \includegraphics[width=\linewidth]{example-image}
    \caption{Zrzut ekranu z fragmentem dokumentacji dla metody sortującej tablicę liczb całkowitych. 
    Pokazana jest nazwa klasy wraz z opisem, kto jest jej autorem.}
    \label{fig:java_sort_alg_screenshot}
\end{figure}

\subsection{Zadanie 10.}

Proszę zweryfikować poprawność otrzymanych i wykonanych implementacji algorytmów poprzez napisanie testów jednostkowych.
\underline{Każda klasa} implementująca algorytm powinna być przetestowana. 

Proszę to zadanie wykonać w sposób przemyślany (kopiowanie i powielanie kodu testów pomiędzy klasami jest ostatnim
krokiem, na który należy się zdecydować).
W celu zniwelowania redundancji kodu testów, proszę wykorzystać cechy paradygmatu programowania obiektowego.

\end{document}
